% Standalone problem document, generated from the adjacent README.md.
% Compile with XeLaTeX. Mathematical content is maintained in Markdown.
\documentclass[11pt,a4paper]{article}
\usepackage[fontsize=10.5pt]{scrextend}
\usepackage[margin=22mm,headheight=15pt,headsep=9mm,footskip=11mm]{geometry}
\usepackage{fontspec}
\setmainfont{texgyrepagella}[Extension=.otf,UprightFont=*-regular,BoldFont=*-bold,ItalicFont=*-italic,BoldItalicFont=*-bolditalic]
\setsansfont{texgyreheros}[Extension=.otf,UprightFont=*-regular,BoldFont=*-bold,ItalicFont=*-italic,BoldItalicFont=*-bolditalic]
\setmonofont{texgyrecursor}[Extension=.otf,UprightFont=*-regular,BoldFont=*-bold,ItalicFont=*-italic,BoldItalicFont=*-bolditalic,Scale=MatchLowercase]
\usepackage{amsmath,amssymb}
\usepackage{unicode-math}
\setmathfont{texgyrepagella-math.otf}
\usepackage{xcolor,microtype,enumitem,needspace,fancyhdr}
\usepackage{bookmark}
\definecolor{ink}{HTML}{17344B}
\definecolor{accent}{HTML}{197D80}
\definecolor{muted}{HTML}{596773}
\hypersetup{colorlinks=true,linkcolor=accent,urlcolor=accent,pdftitle={MF-05 - Local
Hölder continuity of the joint spectral
radius},pdfauthor={Open Problems in Numerical Linear Algebra}}
\urlstyle{same}
\setlength{\parindent}{0pt}
\setlength{\parskip}{5pt plus 1pt minus 1pt}
\linespread{1.04}
\setlength{\emergencystretch}{3em}
\widowpenalty=10000
\clubpenalty=10000
\displaywidowpenalty=10000
\allowdisplaybreaks[1]
\setlist{nosep,leftmargin=1.5em}
\providecommand{\tightlist}{\setlength{\itemsep}{0pt}\setlength{\parskip}{0pt}}
\providecommand{\pandocbounded}[1]{#1}
\pagestyle{fancy}
\fancyhf{}
\fancyhead[L]{\sffamily\footnotesize\color{muted}OPEN PROBLEMS IN NUMERICAL LINEAR ALGEBRA}
\fancyhead[R]{\sffamily\bfseries\footnotesize\color{accent}MF-05}
\fancyfoot[L]{\parbox[b]{0.9\textwidth}{\sffamily\fontsize{8}{10}\selectfont\color{muted}Editorial ratings; a bounded search cannot certify that a problem remains unsolved.\\Literature check: 2026-09-10}}
\fancyfoot[R]{\sffamily\footnotesize\color{muted}\thepage}
\renewcommand{\headrulewidth}{0.3pt}
\renewcommand{\headrule}{\hbox to\headwidth{\color{accent}\leaders\hrule height \headrulewidth\hfill}}
\makeatletter
\renewcommand{\section}{\Needspace{5\baselineskip}\@startsection{section}{1}{0pt}{12pt plus 2pt minus 2pt}{4pt}{\sffamily\bfseries\large\color{ink}}}
\renewcommand{\subsection}{\Needspace{4\baselineskip}\@startsection{subsection}{2}{0pt}{9pt plus 1pt minus 1pt}{3pt}{\sffamily\bfseries\normalsize\color{ink}}}
\makeatother
\setcounter{secnumdepth}{0}
\begin{document}
{\sffamily\small\color{muted}Matrix functions and stability\par}
\vspace{3pt}
{\raggedright\hyphenpenalty=10000\exhyphenpenalty=10000\sffamily\bfseries\fontsize{21}{24}\selectfont\color{ink}Local
Hölder continuity of the joint spectral radius\par}
\vspace{7pt}
{\sffamily\small\textbf{Difficulty:} challenging\quad\textbf{Importance:} interesting
to the community\par}
{\sffamily\small\bfseries\color{accent}Status: Partially resolved\par}
\vspace{4pt}
{\color{accent}\hrule height 0.8pt}
\vspace{6pt}
\textbf{Rating rationale:} Challenging because reducible families
obstruct uniform perturbation estimates; community impact comes from
conditioning and robustness of stability computations.

\subsection{Context and notation}\label{context-and-notation}

Let \(\mathcal H_d\) denote the nonempty compact subsets of
\(\mathbb C^{d\times d}\). Use the spectral norm and its Hausdorff
distance

\[
d_H(\mathcal M,\mathcal N)=\max\left\{
\sup_{A\in\mathcal M}\inf_{B\in\mathcal N}\|A-B\|_2,
\sup_{B\in\mathcal N}\inf_{A\in\mathcal M}\|A-B\|_2\right\}.
\]

The joint spectral radius is

\[
\widehat\rho(\mathcal M)=\lim_{k\to\infty}
\max_{A_1,\ldots,A_k\in\mathcal M}\|A_k\cdots A_1\|_2^{1/k}.
\]

These definitions also apply to finite real matrix sets. The ordinary
spectral radius of one matrix is written \(\rho(A)\).

\subsection{Problem statement}\label{problem-statement}

For every \(d\ge1\) and \(\mathcal M_0\in\mathcal H_d\), do there exist
\(r,C>0\) such that

\[
|\widehat\rho(\mathcal M)-\widehat\rho(\mathcal N)|
\le C d_H(\mathcal M,\mathcal N)^{1/d}
\]

whenever \(d_H(\mathcal M,\mathcal M_0)<r\) and
\(d_H(\mathcal N,\mathcal M_0)<r\)?

\subsection{Reference and status
evidence}\label{reference-and-status-evidence}

Epperlein and Wirth, \href{https://arxiv.org/html/2311.18633v2}{The
joint spectral radius is pointwise Hölder continuous}, Linear Algebra
Appl. 704 (2025), 92--122,
\href{https://doi.org/10.1016/j.laa.2024.09.016}{DOI}, §2, Conjecture 3
(L1). Pointwise results in that paper do not prove this local assertion.

\subsection{Additional status
evidence}\label{additional-status-evidence}

Searches combining ``joint spectral radius'' with ``local Hölder'',
``Lipschitz lower'', ``trajectory bounds'', the authors' names, and
2025/2026 found no later resolution. These searches supplement the
explicit 2025 conjectures; they do not establish exhaustiveness. The
three entries are separately named assertions in the source, not a count
of dimensional cases.

\subsection{Audit --- 2026-09-10}\label{audit-2026-09-10}

Rechecked \href{https://arxiv.org/html/2311.18633v2}{Epperlein--Wirth,
§1 and Conjecture 3 (L1)}. Local Lipschitz continuity settles the
irreducible-family subcase, but the displayed exponent for arbitrary
compact families remains open. Author and local-Hölder follow-up
searches found no resolution; the weaker general pointwise results do
not suffice.
\end{document}
